\part{Développements en langage de commande du shell}
\label{part:dev_sh}

\section{Type de shell}
   Pour les scripts simples reposant sur peu de fonctionnalités
du shell, il est préférable de privilégier la portabilité et de
se limiter à une syntaxe POSIX. On pourra tester, par exemple,
que les scripts fonctionnement avec plusieurs interpréteurs (dash, bash, ...)

   Pour les scripts plus complexes, l'usage actuel est d'utiliser préférentiellement
bash, car il est largement répandu. Quelques scripts existent en ksh.

\smallbreak
  Les codes shell doivent débuter par un sheebang sous la forme suivante : \\
 (ex pour un shell \texttt{bash}): 
\begin{boxedverbatim}
#!/usr/bin/env  bash
\end{boxedverbatim}

\section{Bibliothèque de fonctions}
   Les fonctions les plus courantes ont été regroupées dans \url{AIDA_shell_lib.sh}, 
qui est installé dans  \url{$AIDA_HOME/lib}. 
   Pour l'utiliser, voici un exemple typique de code d'initialisation, à placer
en début de script : 

\begin{lstlisting}[language=sh,frame=trbl,basicstyle=\small\ttfamily]
AIDA_LIB=${AIDA_LIB:-${AIDA_HOME}/lib}
SHELL_LIB=${AIDA_LIB}/AIDA_shell_lib.sh
if ! test -f $SHELL_LIB
then
  echo "Bibliotheque de fonctions shell absente ! ${SHELL_LIB}"
  exit
fi
. ${SHELL_LIB}

AIDA_env_init
export PATH=${AIDA_BIN}:$PATH
\end{lstlisting}

\section{Exemples de code}
\subsection{Travail sur des données particulières}
  Pour offrir la possibilité de travailler sur fichiers (données, exécutables,...) autres
que que ceux de l'arborescence standard, l'utilisateur à la possibilité, par des
variables d'environnement AIDA\_DATA, AIDA\_BIN, ... de spécifier un chemin
particulier (exemple : un cd-rom) pour un type de fichiers.

   Pour prendre cela en compte dans un script, la méthode
la plus adaptée est la suivante :
\begin{verbatim}
AIDA_HOME=${AIDA_HOME}
AIDA_DATA=${AIDA_DATA:${AIDA_HOME}/data}
AIDA_BIN=${AIDA_DATA:${AIDA_HOME}/bin}
\end{verbatim}

   Pour les scripts qui utilisent la bibliothèque \url{AIDA_shell_lib.sh}, 
ces initialisations sont faites dans la fonction \texttt{AIDA\_env\_init}

\subsection{Récupérer des propriétés d'un vecteur dans son en-tête}
Le plus portable est d'utiliser \texttt{vec2hdr} et \texttt{FDDfiltre} 
(les solutions à base de \texttt{strings}, \texttt{grep} peuvent
donner des résultats différents selon les versions Unix/linux).
Par exemple, pour connaître l'unité : 
\begin{verbatim}
 unit_str=`vec2hdr $Acq_id $VECT | FDDfiltre -l -f 3 -c unit`
\end{verbatim}
